\section{Extracción seleccionada y tratamiento hidráulico de bancos (ENV-19)}
\label{sec:sd4-ambiente-env19}

ENV-19 conserva dos bancos independientes, 900 m$^3$/h efectivos por banco y las condiciones de ENV-16. Selecciona ocho Systemair \texttt{EX 180A-4}, artículo \texttt{135279}, variante 400 V trifásica/50 Hz: dos ventiladores en paralelo por tren, dos trenes independientes por banco. Cada tren recibe las cuatro captaciones de su banco; uno trabaja y el otro queda disponible. La pérdida de un ventilador invalida todo su tren y arranca el otro; no se declara que el ventilador restante mantenga 900 m$^3$/h. Cada ramal del ventilador detenido tiene aislamiento antirretorno, y el cambio de tren se verifica mediante caudal efectivo, no por contacto de motor. Se incluyen ocho kits de adaptación \texttt{EX180A-160/135704}; la identificación del kit no acredita las compuertas ni su aptitud para gas.

ENV-22 utiliza el certificado DNV emisión 7 de 2024, con límite de +50 grados C para EX180A-4 y correspondencia placa/motor pendiente; ese límite prevalece sobre el +60 de la ficha histórica siguiente. La ficha de esa variante publica 248 W, 0,73 A, máximo de 869 m$^3$/h, 6,4 kg, ambiente/aire de $-20$ a 60\,$^{\circ}$C y clasificación \texttt{II 2 G Ex db eb h IIB+H2 T4 Gb}, certificado \texttt{Presafe 16ATEX8598X}. La presencia explícita de H$_2$ es necesaria: IIB sin esa extensión no basta. Su curva gráfica a 400 V proporciona aproximadamente 200 Pa a 450 m$^3$/h; la ficha calcula otro punto de 406 m$^3$/h/225 Pa, que no se traslada a 450. La lectura gráfica justifica una preselección con presión de proyecto menor, no un ensayo ni una garantía numérica del fabricante. Se recibe la selección numérica de los dos ventiladores en paralelo, condiciones especiales del certificado y protección PTC autorizada antes de compra \parencite[pp.~1--5 y 7]{lote19-systemair-ex180}.

\textbf{Red de extracción de baja pérdida.} Se sustituyen los dieciséis ramales de 125 mm por 160 mm interiores, con ocho activos; los cuatro colectores de 250 mm pasan a 355 mm interiores. Cada rama de captación conduce 225 m$^3$/h; cada colector, 900; las dos conexiones de motor, 450 cada una a 160 mm. Longitudes máximas de obra: ramal 8 m, colector 15 m y conexión de motor 0,8 m. Con densidad propia conservadora 1,20 kg/m$^3$, factor Darcy 0,025 y coeficientes de accesorios propios $K=5$, 4 y 1 respectivamente, las velocidades son 3,1085, 2,5258 y 6,2170 m/s. La pérdida de la ruta más desfavorable se calcula como
\[
\Delta p=\sum\left(f\frac{L}{D}+K\right)\frac{\rho v^2}{2}
       +25_{\rm filtro}+15_{\rm terminal}+15_{\rm equilibrado}+20_{\rm reserva}
       \simeq157\ {\rm Pa}<180\ {\rm Pa}.
\]
Los 25 Pa de filtro sucio son límite de mantenimiento de proyecto; no dato de un filtro adquirido. El presupuesto incluye la conexión estrecha de cada motor, no sólo el colector. Toda compuerta, sellado, reducción o accesorio real se incorpora al presupuesto; si lo eleva sobre 180 Pa, se amplía ruta o sustituye conjunto, sin bajar caudal. Los filtros/elementos en la corriente de H$_2$ deben ser compatibles con clasificación, descarga electrostática y operación continua. La curva paralela y el ajuste de equilibrio deben excluir oscilación o recirculación entre ventiladores. La aceptación mide 225 m$^3$/h en cada captación y 900 por banco con pérdida de un tren y filtro en límite; no sirve demostrar 900 sólo en descarga sin reparto.

Cuatro motores activos reservan 0,992 kW y $\sqrt{3}(400)(4)(0,73)/1000=2,0230$ kVA; los ocho durante un relevo conservador reservan 1,984 kW/4,0461 kVA antes de arranque y de auxiliares. Cada fase absorbe 2,92 A con cuatro y 5,84 A con ocho, como corriente de placa, no medida. Los 248 W no se usan para deducir potencia aparente ni FP; el control de gas, PTC, compuertas y flujo tiene aporte adicional. La fuente durante batería y el transitorio se resuelven en CEN/ENE antes de habilitar: seleccionar el motor no acredita capacidad eléctrica restante.

\textbf{Tratamiento de reposición por etapas.} El balance neto de 17,5161 kW por banco representa exterior a consigna, pero no la potencia de una batería fría que debe condensar agua. Para 900 m$^3$/h, $\dot m_{da}=900(1,094636)/3600=0,273659$ kg/s; con $w_e=0,0289204$, $w_i=0,0083562$ y $h=1,006T+w(2501+1,86T)$, el aire saturado con humedad de consigna tiene rocío aproximado de 11,35\,$^{\circ}$C. El enfriamiento exterior hasta ese punto exige 21,0532 kW y elimina 20,2593 kg/h; recalentarlo a 24\,$^{\circ}$C demanda 3,5371 kW. La diferencia conserva 17,5161 kW, sin sumar otra vez sensible/latente.

La esquina 23\,$^{\circ}$C/40\,\% HR limita $w$ a aproximadamente 0,006975 kg/kg. Se fija rocío de proyecto de 8,5\,$^{\circ}$C para la reposición cuando se requiera esa esquina: $w=0,006888$, frío 22,8604 kW, agua 21,7058 kg/h y recalentamiento hasta 24\,$^{\circ}$C de 4,3215 kW por banco. Son capacidades mínimas de proceso calculadas; la salida real de serpentín incluye factor de bypass, agua, ventilador y ensuciamiento. No se declara que una batería entregue aire saturado a su temperatura de agua. El control modula frío y recalentamiento según rocío/T/HR; no mantiene siempre 8,5\,$^{\circ}$C y luego humedece sin necesidad. Calor de baterías/cargador, envolvente y aire de impulsión se agregan por estado en la selección real. La extracción alta no se recircula a la unidad de tratamiento.

Se elige arquitectura de agua fría con dos Carrier \texttt{30RB065-SS}, cada uno capaz de atender el tratamiento de ambos bancos cuando su selección corregida lo confirme; uno de servicio y otro de reserva. Cuatro unidades exteriores de tratamiento, dos por banco, se construyen con serpentín frío, separador de gotas, bandeja/drenaje supervisado, recalentamiento y ventilador de impulsión; cada una debe tratar 900 m$^3$/h completos. Las unidades hidráulicas y sus bombas/rutas permanecen independientes hasta los ramales aislables, para que mantenimiento no retire ambos bancos. Las cargas nominales de los cuatro tratamientos no se suman como cuatro bancos, pero el relevo simultáneo y los arranques sí se evalúan.

La ficha Carrier individualiza 63 kW nominales, R410A, 400 V/3 fases/50 Hz, dos compresores, pasos 0/50/100\,\%, caudal nominal de agua 10,8 m$^3$/h, conexión DN50, envolvente 2,240 por 1,150 por 2,195 m y masa en operación 790 kg. Se incorpora opción real \texttt{003A Gold Fin}; la opción no demuestra por sí sola aptitud marina de toda la instalación. El rango exterior publicado 10--48\,$^{\circ}$C cubre 35\,$^{\circ}$C, pero no acredita 63 kW a agua 5/10\,$^{\circ}$C, ni fuera de ese rango. La cifra nominal tampoco incluye curva sensible/latente del serpentín, caudal/presión de impulsión, potencia máxima absorbida o mínimos de ciclo. Antes de liberar compra se requiere selección de fabricante y tratamiento a medida en esos puntos, régimen anual y disipación real por estado \parencite[pp.~6--8 y 11]{lote19-carrier-30rb065}.

El escalón mínimo nominal de 31,5 kW es superior a varias cargas de reposición de clima templado. Se prescribe depósito de inercia aislado de 1.500 L útiles por circuito: con salto útil propio de 3 K almacena $1500(4,18)(3)/3600=5,225$ kWh. Frente a 31,5 kW y carga nula permite aproximadamente 597 s antes de agotar ese salto; no acredita un tiempo mínimo ON ni OFF desconocido. La compra debe conciliar límites del fabricante, caudal mínimo, anticongelación, control a carga baja y disponibilidad durante batería. La fuente de recalentamiento se selecciona en ENV-25: VEAB con control externo apto para 40 grados C, 18 kW máximos de calor y 0,100 kW de controles para dos rutas activas. No se libera una fuente hidráulica ni se supone recuperación de condensación no publicada para esta variante.

\textbf{Implantación y frontera de seguridad.} Se fija para cada chiller un recinto de obra exterior de 5 por 4 m, con altura libre abierta y acceso perimetral; dentro se representa la envolvente de catálogo de 2,240 por 1,150 m. La reserva no sustituye las separaciones del manual de instalación. Se fijan cuatro módulos exteriores de tratamiento de obra de 2 por 1 por 2 m y dos reservas hidráulicas de 2 por 2 m para inercia/bombeo. Esas envolventes son límites de adquisición del conjunto a medida, no dimensiones de catálogo ya verificadas. Cuatro bastidores de extracción de obra de 1,2 por 0,8 m contienen dos motores cada uno; el plano distingue sus huellas de la unidad Carrier. Las descargas altas quedan a distancia horizontal propia de al menos 3 m de tomas de aire y chillers, sin retorno a puerta o admisión. El paso de tubería hidráulica y drenaje conserva separación de tableros, barras y salida de evacuación, con bandeja y detección MON-16. El piso de chiller recibe como mínimo 790 kg más tubería/base y cargas dinámicas calculadas por estructura.

La arquitectura de bancos no lleva refrigerante a sus recintos: el R410A permanece en el circuito hermético exterior y sólo entra agua en los serpentines. Se reciben ficha de seguridad, carga total, detección/ventilación de fugas exterior, accesos y protección contra heladas. Esta decisión no regulariza los dos circuitos R32 de la pareja candidata s-MEXT en TI; su selección y seguridad siguen sujetas al volumen, carga y accesorios reales de ese recinto. RT-06.03 permanece por completar por unidades a medida y holguras/rutas finales; RT-06.13 permanece por completar por mapas térmicos/mínimos, calor de bancos, recalentamiento, fuente y consumo del conjunto, enclavamiento y selección numérica de extracción. La elección de arquitectura y variante de motor reduce esas brechas sin afirmar recepción realizada.
